\section{Hashing, Bloom Filter, Skip List, Treap --- risoluzioni}

\subsection*{\qid{5.1}\hot\ Hashing universale: definizione, esempio, prova}

\begin{domanda}
Definisci una classe di funzioni hash universali; fornisci un esempio e dimostrane
l'universalit\`a. \emph{(16/01/23, 07/09/23, 02/11/23, 10/01/25)}
\end{domanda}

\textbf{Definizione.} Una famiglia $\mathcal{H}\subseteq\{h:U\to[0,m)\}$ \`e \emph{universale} se
per ogni coppia di chiavi distinte $x\ne y$ di $U$:
\[
\boxed{\ \Pr_{h\in\mathcal{H}}\bigl[h(x)=h(y)\bigr]\;\le\;\frac{1}{m}\ }
\]
dove $h$ \`e estratta uniformemente da $\mathcal{H}$. Equivalentemente: il numero di funzioni
della famiglia che fanno collidere $x$ e $y$ \`e $\le|\mathcal{H}|/m$. \kw{La casualit\`a sta
nella scelta della funzione, non nelle chiavi}: la garanzia vale per ogni input, anche avversario.

\textbf{Esempio (prodotto scalare modulo primo).} Sia $m$ \kw{primo} e si scrivano le chiavi in
base $m$: $k=(k_0,k_1,\dots,k_{r-1})$ con $k_i\in[0,m)$. Per ogni vettore
$a=(a_0,\dots,a_{r-1})\in[0,m)^r$ si definisce
\[
h_a(k)\;=\;\Bigl(\sum_{i=0}^{r-1}a_i\,k_i\Bigr)\bmod m ,
\]
con $\mathcal{H}=\{h_a\}$, $|\mathcal{H}|=m^r$ (scegliere $h$ a caso $=$ scegliere $a$ a caso).

\textbf{Prova di universalit\`a.} Siano $x\ne y$; differiscono in almeno una cifra, diciamo la
$d$-esima: $x_d\ne y_d$. La collisione $h_a(x)=h_a(y)$ equivale a
\[
a_d\,(x_d-y_d)\;\equiv\;-\!\!\sum_{i\ne d}a_i\,(x_i-y_i) \pmod m .
\]
Poich\'e $m$ \`e primo e $x_d-y_d\not\equiv0$, \kw{$(x_d-y_d)$ \`e invertibile in $\mathbb{Z}_m$}:
fissati comunque gli altri $r-1$ coefficienti, esiste \kw{esattamente un} valore di $a_d$ che
soddisfa l'equazione:
\[
a_d\;\equiv\;-(x_d-y_d)^{-1}\!\!\sum_{i\ne d}a_i\,(x_i-y_i)\pmod m .
\]
Le funzioni che fanno collidere $x,y$ sono dunque $m^{r-1}$, e
$\Pr[h_a(x)=h_a(y)]=\frac{m^{r-1}}{m^{r}}=\frac{1}{m}$: universale. $\blacksquare$

\subsection*{\qid{5.2} Universale $+$ chaining: catene attese $\OO{1}$}

\begin{domanda}
Definisci una classe universale e mostra che, usata nell'hashing con chaining, ottiene lunghezza
attesa $\OO{1}$ per catena quando $m=\OO{n}$; da qui il calcolo del numero atteso di collisioni.
\emph{(06/11/23; esercizi collegati in 10/07/20 e 11/11/20)}
\end{domanda}

Definizione di classe universale: come \qid{5.1}. Sia $T[0,m-1]$ una tabella con chaining, $h$
estratta da una famiglia universale, $n$ chiavi inserite. Fissata una chiave $x$, sia $C_x$ il
numero di chiavi $y\ne x$ che collidono con $x$. \kw{Per linearit\`a dell'attesa e
universalit\`a}:
\[
\mathbb{E}[C_x]\;=\;\sum_{y\ne x}\Pr[h(y)=h(x)]\;\le\;\frac{n-1}{m}\;<\;\alpha:=\frac{n}{m}.
\]
Quindi il costo atteso di una ricerca \`e $\boxed{\OO{1+\alpha}}$, che \`e $\OO{1}$ non appena
$m=\Omega(n)$: \kw{$m=\OO{n}$ celle bastano per catene di lunghezza attesa costante}.
$\blacksquare$

\textbf{Esercizio collegato (10/07/20, 11/11/20): collisioni totali attese.} Il numero di coppie
collidenti \`e $C=\sum_{\{x,y\}}\mathbf{1}[h(x)=h(y)]$ su $\binom{n}{2}$ coppie:
\[
\mathbb{E}[C]\;\le\;\binom{n}{2}\frac{1}{m}\;=\;\frac{n(n-1)}{2m}.
\]
Con i numeri d'esame ($n=10$, range $[0,22]\Rightarrow m=23$):
$\mathbb{E}[C]\le\frac{10\cdot9}{2\cdot23}=\frac{45}{23}\approx1.96$.

\subsection*{\qid{5.3}\hot\ Cuckoo hashing: teorema dei cammini}

\begin{domanda}
Enuncia e dimostra il teorema sul cuckoo graph: per ogni coppia di posizioni $i,j$ della tabella
cuckoo, la probabilit\`a che il cammino minimo $i\to j$ abbia lunghezza $L$ \`e al pi\`u
$c^{-L}/r$, dove $r$ \`e il numero di posizioni della tabella, $n$ le chiavi memorizzate e
$r\ge 2cn$. Commenta l'impatto sul load factor. \emph{(10/01/20, 10/02/23, 03/06/25)}
\end{domanda}

\textbf{Contesto.} Nel cuckoo hashing ogni chiave $k$ ha due posizioni possibili
$h_1(k),h_2(k)$ in una tabella di $r$ celle; l'inserimento sfratta a catena. Il \emph{cuckoo
graph} ha come nodi le $r$ celle e, per ogni chiave, un arco fra le sue due posizioni. Le
prestazioni dipendono dall'assenza di cammini lunghi e cicli.

\textbf{Teorema.} Se le funzioni hash sono (idealmente) casuali e $r\ge 2cn$ per una costante
$c>1$, allora per ogni coppia di celle $i,j$ e ogni $L\ge1$:
\[
\boxed{\ \Pr\bigl[\exists\text{ cammino minimo }i\to j\text{ di lunghezza }L\bigr]
\;\le\;\frac{c^{-L}}{r}\ }
\]

\textbf{Prova (induzione su $L$).}

\emph{Base $L=1$.} Esiste un arco $(i,j)$ sse qualche chiave $k$ ha $\{h_1(k),h_2(k)\}=\{i,j\}$.
Per una singola chiave: $\Pr=\frac{2}{r^2}$ (due orientamenti). Per union bound sulle $n$ chiavi,
\kw{usando l'ipotesi $r\ge2cn$, cio\`e $n\le r/(2c)$}:
\[
\Pr[\exists\,\text{arco } i\text{--}j]\;\le\;\frac{2n}{r^{2}}
\;\le\;\frac{2}{r^{2}}\cdot\frac{r}{2c}\;=\;\frac{c^{-1}}{r}.
\]

\emph{Passo.} Un cammino minimo $i\to j$ di lunghezza $L$ \`e composto da un cammino minimo di
lunghezza $L-1$ da $i$ a una cella intermedia $\ell$, seguito da un arco $(\ell,j)$. Per union
bound su $\ell$, con ipotesi induttiva e indipendenza dell'ultimo arco:
\[
\Pr[\text{cammino }i\to j\text{ di lung. }L]
\;\le\;\sum_{\ell=1}^{r}\frac{c^{-(L-1)}}{r}\cdot\frac{c^{-1}}{r}
\;=\;r\cdot\frac{c^{-(L-1)}}{r}\cdot\frac{1}{c\,r}
\;=\;\frac{c^{-L}}{r}. \;\blacksquare
\]

\textbf{Conseguenze e load factor.} Probabilit\`a che $i$ e $j$ siano connessi (lunghezza
qualunque): $\sum_{L\ge1}\frac{c^{-L}}{r}=\frac{1}{r(c-1)}$. Da qui:
(i) \kw{insert atteso $\OO{1}$} (cammini attesi di lunghezza costante) e \kw{ricerca/delete
$\OO{1}$ worst case} (si guardano 2 celle);
(ii) la probabilit\`a di un \emph{ciclo} (fallimento $\Rightarrow$ rehash) si maggiora con
$\frac{1}{c-1}$: con $c=3$ \`e $\le\frac12$, quindi $\OO{1}$ rehash attesi, ammortizzato $\OO{1}$.
\emph{Prezzo}: $r\ge2cn=6n$ celle, cio\`e \kw{load factor $\le1/6\approx17\%$}: il cuckoo classico
garantisce i bound solo a tabelle molto scariche (varianti con pi\`u hash o celle multi-slot
arrivano oltre il 90\%).

\subsection*{\qid{5.4}--\qid{5.5}\hot\ Bloom Filter: definizione, bit a 0, falso positivo, ottimo}

\begin{domanda}
\qid{5.4} Calcola la probabilit\`a che una posizione del bit-array del Bloom Filter (taglia $m$,
$n$ chiavi inserite) sia 0, e deriva/dimostra la probabilit\`a d'errore complessiva. Variante:
dimostra che tale probabilit\`a \`e $\tfrac12$ quando il numero di funzioni hash \`e quello
ottimo. \emph{(11/11/20, 27/05/24, 11/11/24, 15/07/25)}\\
\qid{5.5} Definisci un Bloom Filter su un dizionario di $n$ chiavi e le operazioni supportate;
enuncia e dimostra il risultato sul falso positivo. \emph{(24/07/23)}
\end{domanda}

\textbf{Definizione e operazioni (\qid{5.5}).} Un Bloom Filter per un dizionario $D$ di $n$ chiavi
\`e un array binario $B[0,m-1]$, inizialmente tutto 0, pi\`u $k$ funzioni hash indipendenti
$h_1,\dots,h_k:U\to[0,m)$.
\emph{Insert}$(x)$: pone $B[h_i(x)]=1$ per $i=1..k$; costo $\OO{k}$.
\emph{Query}$(y)$: risponde ``presente'' sse $B[h_i(y)]=1$ \emph{per tutti} gli $i$; costo
$\OO{k}$. Non supporta la cancellazione (creerebbe falsi negativi; la variante counting/spectral
usa contatori). Propriet\`a fondamentale: \kw{niente falsi negativi, possibili falsi positivi}.

\textbf{Probabilit\`a che un bit sia 0.} Gli inserimenti accendono $kn$ posizioni casuali
indipendenti; una posizione fissata resta 0 se tutte le $kn$ estrazioni la evitano:
\[
\boxed{\ p_0=\Bigl(1-\frac{1}{m}\Bigr)^{kn}\;\approx\;e^{-kn/m}\ }
\]

\textbf{Probabilit\`a di falso positivo.} $y\notin D$ \`e dichiarata presente sse le sue $k$
posizioni sono tutte a 1 (assumendo indipendenza, approssimazione standard):
\[
\boxed{\ \varepsilon\;=\;(1-p_0)^{k}\;\approx\;\bigl(1-e^{-kn/m}\bigr)^{k}\ }
\]

\textbf{Numero ottimo di hash e $p_0=\tfrac12$ (\qid{5.4}).} Minimizziamo
$\varepsilon(k)=(1-e^{-kn/m})^k$. Posto \kw{$t=e^{-kn/m}$}, cio\`e $k=-\frac{m}{n}\ln t$:
\[
\ln\varepsilon=k\ln(1-t)=-\frac{m}{n}\,\ln t\,\ln(1-t).
\]
La funzione $g(t)=\ln t\,\ln(1-t)$ \`e \kw{simmetrica rispetto a $t\mapsto1-t$} e unimodale: il
punto ottimo \`e $t=\tfrac12$ (formalmente $g'(t)=\frac{\ln(1-t)}{t}-\frac{\ln t}{1-t}=0$ in
$t=\frac12$). Quindi
\[
e^{-k^{*}n/m}=\tfrac12
\;\Longrightarrow\;
\boxed{\;k^{*}=\frac{m}{n}\ln2\,,\qquad p_0=\tfrac12\;}
\]
\kw{con $k$ ottimo met\`a dei bit \`e a 0 e met\`a a 1} (massima ``entropia'' del filtro).
Errore all'ottimo:
\[
\boxed{\ \varepsilon^{*}=\Bigl(\tfrac12\Bigr)^{k^{*}}=2^{-\frac{m}{n}\ln2}
\approx(0.6185)^{\,m/n}\ }
\]
\kw{esponenzialmente piccolo nel numero di bit per chiave $m/n$}
(es.\ $m/n=10\Rightarrow\varepsilon\approx0.8\%$ con $k^*\approx7$).

\subsection*{\qid{5.6} Intersezione di due insiemi su due server con un Bloom Filter}

\begin{domanda}
Dati due insiemi $A$, $B$ memorizzati su due server, mostra come calcolarne l'intersezione usando
un Bloom Filter e un solo round di comunicazione, ammettendo errori. \emph{(05/06/23)}
\end{domanda}

\emph{Protocollo.} Il server 1 costruisce un Bloom Filter $BF_A$ di $A$ (taglia $m=\Theta(|A|)$
bit, $k$ ottimo) e \kw{lo invia al server 2} (un solo messaggio). Il server 2 esegue
$\textit{Query}_{BF_A}(b)$ per ogni $b\in B$ e restituisce
$X=\{b\in B:\ BF_A(b)=\text{presente}\}$.

\emph{Correttezza/errore.} \kw{Niente falsi negativi $\Rightarrow$ $A\cap B\subseteq X$}: nessun
elemento dell'intersezione \`e perso. Ogni $b\in B\setminus A$ finisce in $X$ con probabilit\`a
$\varepsilon\approx(0.6185)^{m/|A|}$: attesi $\varepsilon\,|B\setminus A|$ falsi positivi,
controllabili aumentando $m$. \emph{Comunicazione}: $m=\OO{|A|}$ \kw{bit} invece di
$\Theta(|A|\log u)$ bit per spedire $A$ esplicitamente. (Se serve l'intersezione esatta, un
secondo round di verifica elimina gli errori --- ma la domanda chiede la versione a un round.)

\subsection*{\qid{5.7} Skip List: altezza media $\OO{\log n}$}

\begin{domanda}
Dimostra che l'altezza (profondit\`a) media di una Skip List con $n$ nodi \`e $\OO{\log n}$.
\emph{(11/11/24, 04/02/25)}
\end{domanda}

\textbf{Setup.} Ogni elemento viene promosso di livello con probabilit\`a $\tfrac12$
indipendentemente: $\Pr[\text{liv}(x)\ge\ell]=2^{-(\ell-1)}$ (geometrica). L'\emph{altezza} \`e
$H=\max_x \text{liv}(x)$.

\textbf{Prova.} Per union bound, per ogni $\ell\ge1$:
\[
\Pr[H\ge\ell]\;\le\;\sum_{x}\Pr[\text{liv}(x)\ge\ell]\;=\;n\,2^{-(\ell-1)}
\qquad(\text{e ovviamente }\le1).
\]
Con l'identit\`a \kw{$\mathbb{E}[H]=\sum_{\ell\ge1}\Pr[H\ge\ell]$}, spezzando la somma in
$\ell\le1+\log_2 n$ e $\ell>1+\log_2 n$:
\[
\mathbb{E}[H]\;\le\;\underbrace{\sum_{\ell=1}^{1+\log_2 n}1}_{\le\,1+\log_2 n}
\;+\;\sum_{j\ge1} n\,2^{-(\log_2 n+j)}
\;=\;1+\log_2 n+\sum_{j\ge1}2^{-j}
\;=\;\boxed{\ \log_2 n+2\ }\;=\;\OO{\log n}. \;\blacksquare
\]
Concentrazione: $\Pr[H\ge c\log_2 n]\le 2n^{1-c}$, cio\`e \kw{altezza $\OO{\log n}$ con alta
probabilit\`a}.

\textbf{Complemento utile (ricerca $\OO{\log n}$ attesa).} Analizzando il cammino di ricerca
\emph{a ritroso}: da ogni nodo si sale di livello con probabilit\`a $\tfrac12$, altrimenti ci si
sposta orizzontalmente; per salire gli $\OO{\log n}$ livelli servono in attesa 2 passi per
livello: $\OO{\log n}$ passi totali.

\subsection*{\qid{5.8}\hot\ Treap: definizione, operazioni, split}

\begin{domanda}
Definisci la struttura Treap su $n$ coppie $\langle$chiave, priorit\`a$\rangle$ con priorit\`a
minima nella radice; descrivi le operazioni supportate e come si implementa lo split su una chiave
$K$ non presente. \emph{(10/02/23, 23/06/25, orale 25/07/22)}
\end{domanda}

\textbf{Definizione.} Un Treap su $n$ coppie $\langle k,\pi\rangle$ \`e un albero binario che \`e
\kw{contemporaneamente}: (i) un \kw{BST rispetto alle chiavi} $k$ (visita in order $=$ ordine
crescente delle chiavi); (ii) uno \kw{heap rispetto alle priorit\`a} $\pi$ --- qui un MIN-heap:
la radice ha la priorit\`a minima. \kw{Se le priorit\`a sono distinte, il Treap \`e unico}: la
radice \`e forzata (priorit\`a minima), le chiavi minori vanno a sinistra, le maggiori a destra,
ricorsivamente. Operazioni: search ($\OO{h}$ come BST), insert e delete tramite rotazioni che
ripristinano lo heap ($\OO{h}$), split e merge ($\OO{h}$); $h=\OO{\log n}$ attesa con priorit\`a
casuali (\qid{5.9}).

\textbf{Split su chiave $K$ non presente.} Obiettivo: dividere $T$ in $T_{<K}$ e $T_{>K}$.
\begin{enumerate}[leftmargin=1.8em,itemsep=1pt]
\item \kw{Si inserisce il nodo fittizio $\langle K,-\infty\rangle$} con l'insert standard: discesa
BST fino a una foglia, poi --- avendo priorit\`a minima di tutti --- il nodo \kw{risale con
rotazioni fino a diventare la radice}.
\item Per la propriet\`a BST, il sottoalbero sinistro della nuova radice contiene esattamente le
chiavi $<K$ e il destro le chiavi $>K$ (le rotazioni preservano BST e heap nei sottoalberi).
\item Si rimuove il nodo fittizio: i suoi due figli sono $T_{<K}$ e $T_{>K}$.
\end{enumerate}
Costo: un insert $=\OO{\log n}$ atteso. (Il \emph{merge} \`e l'operazione inversa: radice fittizia
sopra i due treap, che poi si fa scendere ed eliminare.)

\subsection*{\qid{5.9} Treap: profondit\`a media $\OO{\log n}$}

\begin{domanda}
Enuncia e dimostra che la profondit\`a media di un Treap con priorit\`a scelte uniformemente a
caso \`e logaritmica nel numero di chiavi. \emph{(07/02/20, orale 25/07/22)}
\end{domanda}

\textbf{Teorema.} Se le priorit\`a sono uniformi e indipendenti, per ogni nodo
$\mathbb{E}[\text{depth}(u_j)]\le 2\ln n=\OO{\log n}$.

\textbf{Lemma (struttura degli antenati).} Siano $u_1,\dots,u_n$ i nodi in ordine di chiave.
Allora \kw{$u_i$ \`e antenato di $u_j$ $\iff$ $u_i$ ha la priorit\`a minima nell'intervallo di
chiavi} $S_{ij}=\{u_{\min(i,j)},\dots,u_{\max(i,j)}\}$.

\emph{Prova del lemma.} Sia $z$ il nodo di priorit\`a minima in $S_{ij}$. Nel (unico) Treap, il
pi\`u piccolo sottoalbero che contiene tutto $S_{ij}$ ha come radice un nodo di $S_{ij}$ (le
chiavi di $S_{ij}$ sono contigue, e la radice che le separa ha chiave interna all'intervallo);
per la propriet\`a di heap tale radice \`e $z$, che quindi \`e antenato di tutti i nodi di
$S_{ij}$. Se $z=u_i$: $u_i$ antenato di $u_j$. Se $z\ne u_i$: o $z=u_j$ (e allora \`e $u_j$
l'antenato), oppure $u_i$ e $u_j$ stanno in sottoalberi \emph{diversi} di $z$ (chiavi da parti
opposte della chiave di $z$) e $u_i$ non \`e antenato di $u_j$. In tutti i casi: antenato $\iff$
minimo. $\square$

\emph{Prova del teorema.} Con priorit\`a uniformi, ogni nodo di $S_{ij}$ ($|i-j|+1$ elementi) \`e
il minimo con la stessa probabilit\`a:
\[
\Pr[u_i\text{ antenato di }u_j]=\frac{1}{|i-j|+1}.
\]
La profondit\`a di $u_j$ \`e il numero dei suoi antenati; \kw{per linearit\`a dell'attesa}:
\[
\mathbb{E}[\text{depth}(u_j)]
=\sum_{i\ne j}\frac{1}{|i-j|+1}
=\sum_{t=2}^{j}\frac{1}{t}+\sum_{t=2}^{\,n-j+1}\frac{1}{t}
\;\le\;2(H_n-1)\;\le\;\boxed{\ 2\ln n\ }=\OO{\log n}. \;\blacksquare
\]

\textbf{Commento.} \kw{Il Treap con priorit\`a casuali \`e un random BST}, qualunque sia l'ordine
di inserimento delle chiavi: \`e la sua ragion d'essere.

% ==================================================================
\section{Prefix Search --- risoluzioni}

\subsection*{\qid{6.1} Ricerca lessicografica in un Patricia Trie}

\begin{domanda}
Dato un Patricia Trie su un dizionario $S$ di $n$ stringhe, descrivi come cercare la posizione
lessicografica di $P[1,p]$ in $S$; enuncia la complessit\`a in tempo e in I/O. \emph{(24/07/23)}
\end{domanda}

\textbf{Struttura (richiamo).} Il Patricia Trie \`e il compacted trie in cui ogni arco memorizza
soltanto il \emph{primo carattere} dell'etichetta, e ogni nodo la lunghezza (skip) del prefisso
raggiunto; le stringhe complete risiedono solo nelle foglie (tipicamente su disco).

\textbf{Ricerca (blind search), tre passi:}
\begin{enumerate}[leftmargin=1.8em,itemsep=1pt]
\item \emph{Discesa cieca.} Dalla radice, a ogni nodo con skip $\ell$ si usa $P[\ell+1]$ per
scegliere l'arco (confrontando solo il carattere di diramazione, \kw{senza verificare i caratteri
saltati}). Si termina a una foglia candidata $f$ (in caso di mismatch di diramazione si prosegue
comunque su un rappresentante del sottoalbero).
\item \emph{Verifica.} Si carica \kw{l'unica stringa} $s_f$ della foglia e si calcola
$\ell^{*}=\mathrm{lcp}(P,s_f)$. Propriet\`a chiave che rende corretta la blind search:
\kw{$s_f$ massimizza l'lcp con $P$ fra tutte le stringhe di $S$} (le stringhe con lcp maggiore
starebbero sullo stesso cammino di discesa).
\item \emph{Riposizionamento.} Si risale al punto di profondit\`a $\ell^{*}$ del cammino;
confrontando $P[\ell^{*}+1]$ con i caratteri di diramazione si determina la posizione di $P$ fra i
sottoalberi, e contando le foglie a sinistra la sua posizione lessicografica in $S$.
\end{enumerate}

\textbf{Complessit\`a.}
\[
\boxed{\ \OO{p}\ \text{tempo},\qquad \OO{1}\ \text{I/O}\ }
\]
(pi\`u precisamente $\OO{p/B}$ I/O per leggere la stringa candidata). \kw{Una sola stringa
toccata per qualunque ricerca}: \`e il punto di forza del Patricia.

\subsection*{\qid{6.2}--\qid{6.3} Compacted Trie vs Patricia Trie, e l'indice two-level}

\begin{domanda}
\qid{6.2} Descrivi la differenza fra Compacted Trie e Patricia Trie in termini di spazio e tempo
di query, in funzione di $n$, della lunghezza totale $N$ e della lunghezza del pattern $P$.
\emph{(06/03/25)}\\
\qid{6.3} Qual \`e la differenza principale fra Compacted Trie e Patricia Trie, e come \`e stata
sfruttata per costruire un indice two-level efficiente per un insieme di stringhe?
\emph{(03/06/25)}
\end{domanda}

\textbf{Compacted trie.} Archi etichettati con \emph{sottostringhe intere} (coppie
(puntatore, lunghezza) nelle stringhe di $S$). Nodi $\OO{n}$, ma per navigare servono le stringhe:
spazio complessivo $\TT{N}$ caratteri $+$ $\OO{n}$ puntatori. Query: confronta \emph{tutti} i
caratteri di $P$ durante la discesa: tempo $\OO{p}$, ma se le stringhe sono su disco ogni arco
pu\`o costare accessi (pi\`u di $\OO{1}$ I/O).

\textbf{Patricia trie.} Un carattere $+$ skip per arco: \kw{spazio $\OO{n}$ parole, indipendente
da $N$}; query $\OO{p}$ con \kw{un solo accesso a una stringa} per la verifica finale: $\OO{1}$
I/O (\qid{6.1}). Prezzo: i caratteri saltati non sono controllati in discesa (serve la verifica).

\textbf{Differenza principale (\qid{6.3}).} Il compacted trie richiede le etichette complete
(spazio $\propto N$); \kw{il Patricia riduce ogni etichetta a 1 carattere $+$ 1 intero: spazio
indipendente dalla lunghezza delle stringhe}, al costo di una verifica posticipata.

\textbf{Sfruttamento nel two-level index.} Dizionario ordinato su disco a blocchi; in memoria un
\kw{Patricia trie su una stringa campione per blocco} (es.\ la prima):
\begin{itemize}[leftmargin=1.6em,itemsep=1pt]
\item lo spazio in memoria \`e $\OO{n/B_{\mathrm{blk}}}$ parole e \kw{non dipende dalla lunghezza
delle stringhe campione} --- la propriet\`a sfruttata \`e esattamente questa;
\item una query fa blind search in $\OO{p}$ tempo e $\OO{1}$ I/O, individuando \kw{l'unico blocco}
che pu\`o contenere $P$;
\item si legge quel blocco ($\OO{1}$ I/O $+$ $\OO{occ/B}$ per output lunghi) e lo si scandisce;
i blocchi sono compressi con \emph{front coding}, decodificabile dal capo-blocco memorizzato per
intero.
\end{itemize}
Risultato: prefix search in $\OO{p}$ tempo e $\OO{1+occ/B}$ I/O con memoria interna minima.

\subsection*{\qid{6.4} Storage di un dizionario di stringhe: tre soluzioni a confronto}

\begin{domanda}
Dato un dizionario $D$ di $n$ stringhe di lunghezza variabile e lunghezza totale $N$, discuti
almeno 3 soluzioni di memorizzazione, commentando spazio e costo tempo/I/O di
$\texttt{Access}(i)$. \emph{(27/05/24)}
\end{domanda}

\textbf{(1) Concatenazione $+$ array di offset.} Stringhe concatenate in un array di $N$
caratteri; array $\mathit{off}[1,n]$ ($n\lceil\log_2 N\rceil$ bit).
\emph{Access}$(i)$: lettura diretta di $[\mathit{off}[i],\mathit{off}[i+1])$:
\kw{tempo $\OO{|s_i|}$, I/O $\OO{1+|s_i|/B}$} --- ottimo. Spazio: $N$ caratteri $+$ $n\log N$ bit
(gli offset possono dominare con stringhe corte). Nessuna compressione.

\textbf{(2) Concatenazione con lunghezze in linea (senza puntatori).} Ogni stringa preceduta dalla
propria lunghezza (es.\ in $\gamma$-code): spazio $\approx N+\OO{n\log\frac{N}{n}}$ bit.
\emph{Access}$(i)$ richiederebbe la scansione dall'inizio: si corregge \kw{campionando un offset
esplicito ogni $k$ stringhe} (spazio extra $\frac{n}{k}\log N$ bit): Access $=$ salto al campione
$+$ scansione di $\le k$ stringhe: tempo $\OO{k\,\ell_{\max}}$, I/O $\OO{1+k\ell_{\max}/B}$.
\kw{Trade-off spazio/tempo regolato da $k$}.

\textbf{(3) Front coding a blocchi (FC$_B$).} Dizionario ordinato: stringhe consecutive
condividono prefissi lunghi. Blocchi di $k$ stringhe: la prima \kw{per intero} (copia di
ripartenza), le successive come coppie $\langle\ell_j,\hat{s}_j\rangle$ (lcp con la precedente,
suffisso residuo). Spazio: $FC_B(D)\ll N$ su dizionari reali ($\approx-70\%$ sugli URL) $+$
$\frac{n}{k}\log N$ bit di puntatori. \emph{Access}$(i)$: salto al blocco $\lceil i/k\rceil$
($\OO{1}$ I/O) e decodifica in scansione: tempo $\OO{k\,\ell_{\max}}$, I/O tipicamente $\OO{1}$.
\kw{\`E la soluzione usata nei blocchi dell'indice two-level} di \qid{6.3}.

\emph{Confronto}: (1) accesso ottimo, spazio massimo; (2) spazio minore, accesso pi\`u lento
(regolabile); (3) \kw{compressione vera con accesso quasi-ottimo: il compromesso standard}.
(Quarta opzione citabile: compacted trie/Patricia $+$ stringhe su disco, se servono anche ricerche
per prefisso.)

\subsection*{\qid{6.5} Ordinare stringhe con un compacted trie a fan-out hash}

\begin{domanda}
Descrivi come usare un compacted trie (con fan-out dei nodi in hash table) per ordinare $n$
stringhe di lunghezza totale $L$; commenta tempo e spazio. \emph{(10/07/20)}
\end{domanda}

\textbf{Costruzione.} Si inseriscono le $n$ stringhe una alla volta nel compacted trie, dove il
fan-out di ogni nodo \`e una \kw{hash table} (carattere di diramazione $\to$ figlio).
L'inserimento di $s$ scende carattere per carattere: ogni passo costa $\OO{1}$ atteso (lookup
hash); in caso di mismatch a met\`a arco si spezza l'arco con $\OO{1}$ nodi nuovi. Costruzione:
\kw{$\OO{L}$ atteso}; spazio $\OO{n}$ nodi $+$ etichette come puntatori nelle stringhe originali.

\textbf{Ordinamento.} L'ordine lessicografico \`e la \kw{DFS del trie con i figli visitati in
ordine alfabetico}. Ma \kw{le hash table non mantengono l'ordine}: prima della DFS si ordinano le
liste dei figli di ciascun nodo: $\sum_v \OO{d_v\log d_v}\le\OO{n\log\sigma}$ (poich\'e
$\sum_v d_v\le 2n$ e $d_v\le\sigma$). DFS finale: $\OO{n}$.

\textbf{Bilancio.}
\[
\boxed{\ \OO{L+n\log\sigma}\ \text{tempo atteso},\qquad \OO{n}\ \text{spazio extra}\ }
\]
Confronto col lower bound \qid{2.2} $\Om{d+n\log n}$: il trie evita di riconfrontare i prefissi
comuni e sposta il termine di ordinamento a $n\log\sigma$; il prezzo \`e il bound solo
\emph{atteso} (hashing) e la localit\`a peggiore rispetto al Multikey QuickSort.

% ==================================================================
\section{Substring Search --- risoluzioni}

\subsection*{\qid{7.1} Definizioni: Suffix Array e array LCP}

\begin{domanda}
Definisci formalmente il suffix array $SA$ di un testo $T[1,n]$ e il corrispondente array LCP.
\emph{(16/01/23, 02/11/23)}
\end{domanda}

\textbf{Suffix Array.} Dato $T[1,n]$ (terminato da $\$$, minimo lessicografico), $\mathit{SA}[1,n]$
\`e \kw{la permutazione di $\{1,\dots,n\}$ che ordina lessicograficamente i suffissi}:
\[
T[\mathit{SA}[1],n]\;<\;T[\mathit{SA}[2],n]\;<\;\dots\;<\;T[\mathit{SA}[n],n],
\]
cio\`e $\mathit{SA}[i]$ \`e la posizione iniziale dell'$i$-esimo suffisso in ordine
lessicografico. Spazio: $n$ interi $=\Theta(n\log n)$ bit (oltre a $T$).

\textbf{Array LCP.} $\mathit{LCP}[1,n-1]$ dove
\[
\mathit{LCP}[i]\;=\;\bigl|\mathrm{lcp}\bigl(T[\mathit{SA}[i],n],\;T[\mathit{SA}[i+1],n]\bigr)\bigr|
\]
\`e la lunghezza del pi\`u lungo prefisso comune fra i suffissi \kw{lessicograficamente
consecutivi}. Propriet\`a utile: l'lcp fra i suffissi in posizioni $i<j$ \`e
$\min_{i\le t<j}\mathit{LCP}[t]$ (range minimum).

\subsection*{\qid{7.2}, \qid{7.4}, \qid{7.5}\hot\ Ricerca e COUNT con il Suffix Array}

\begin{domanda}
\qid{7.2} Definisci $SA$ e LCP; descrivi come cercare $P[1,p]$ in $SA$; enuncia e dimostra la
complessit\`a della ricerca. \emph{(07/09/23)}\\
\qid{7.4} Mostra come CONTARE in $T[1,n]$ tutte le occorrenze di $P[1,p]$ con un Suffix Array in
memoria; enuncia e dimostra la complessit\`a; qual \`e il costo I/O del RETRIEVAL delle posizioni
se $SA$ \`e su disco? \emph{(10/02/21, 15/01/24)}\\
\qid{7.5} Descrivi come cercare $P[1,p]$ in un $SA$ costruito su $T[1,n]$ e commenta la
complessit\`a in funzione di $p$ e $n$. \emph{(15/07/25)}
\end{domanda}

\textbf{Osservazione chiave.} Le occorrenze di $P$ sono i suffissi che hanno $P$ come prefisso;
nell'ordine lessicografico questi suffissi sono \kw{consecutivi}: formano un intervallo
$\mathit{SA}[i^{*},j^{*}]$.

\textbf{Algoritmo.} \kw{Due ricerche binarie} su $\mathit{SA}$: una per il bordo sinistro $i^{*}$,
una per il bordo destro $j^{*}$. Ogni passo confronta $P$ col suffisso corrente: $\le p$ confronti
di caratteri.

\textbf{Complessit\`a (con prova).}
\[
\boxed{\ \OO{p\log n}\ \text{(ricerca ingenua)},\qquad \OO{p+\log n}\ \text{con Llcp/Rlcp}\ }
\]
\emph{Prova del primo bound}: la ricerca binaria fa $\lceil\log_2 n\rceil$ iterazioni; ogni
iterazione \`e un confronto lessicografico che termina al primo mismatch, dopo al pi\`u $p$
caratteri. \emph{Raffinamento}: mantenendo per l'intervallo corrente $[L,R]$ i valori
$\ell=\mathrm{lcp}(P,\text{suff}_{SA[L]})$, $r=\mathrm{lcp}(P,\text{suff}_{SA[R]})$ e usando gli
array precalcolati $\mathit{Llcp}/\mathit{Rlcp}$, a ogni passo \kw{o si dimezza l'intervallo senza
leggere caratteri, o ogni carattere letto di $P$ fa crescere $\max(\ell,r)$} ($\le p$ volte in
tutto): totale $\OO{p+\log n}$. $\blacksquare$

\textbf{COUNT (\qid{7.4}).} Trovato l'intervallo:
$\boxed{\mathit{occ}=j^{*}-i^{*}+1}$, senza lavoro aggiuntivo. Costo: quello della ricerca.

\textbf{RETRIEVAL con $SA$ su disco (\qid{7.4}).} Le posizioni sono
$\mathit{SA}[i^{*}..j^{*}]$: \kw{celle contigue su disco}, lette con scansione sequenziale in
\[
\boxed{\ \OO{\frac{\mathit{occ}}{B}}\ \text{I/O}\ }
\]
(pi\`u $\OO{\frac{p}{B}\log n}$ I/O per le binarie se anche $T$ \`e su disco). \kw{La contiguit\`a
rende il retrieval I/O-ottimale, output-sensitive.}

\subsection*{\qid{7.3} Costruzione del Suffix Array via qsort}

\begin{domanda}
Definisci $SA$; scrivi lo pseudocodice per costruirlo via qsort; enuncia e commenta la
complessit\`a worst-case (tempo e I/O) di questa costruzione, e il caso in cui il massimo LCP fra
i suffissi \`e limitato da una costante $c$. \emph{(07/09/21, 24/07/23)}
\end{domanda}

\textbf{Pseudocodice.}
\begin{verbatim}
BuildSA-qsort(T[1,n]):
  SA = [1, 2, ..., n]                      // indici dei suffissi
  qsort(SA, cmp)                           // comparatore:
     cmp(i, j) = confronto lessicografico  //   strncmp(T+i, T+j)
                 di T[i,n] con T[j,n]      //   carattere per carattere
  return SA
\end{verbatim}

\textbf{Tempo (worst case).} qsort fa $\OO{n\log n}$ confronti; ogni confronto costa fino a
$\Theta(n)$ caratteri (suffissi con lcp lungo, caso estremo $T=a^n$):
\[
\boxed{\ \OO{n^2\log n}\ \text{tempo worst case}\ }
\]

\textbf{I/O e perch\'e \`e inefficiente.} Due ragioni strutturali:
(a) ogni confronto scandisce fino a $\Theta(n)$ caratteri \kw{a partire da due posizioni
arbitrarie del testo}: $\OO{n/B}$ I/O a confronto, $\OO{\frac{n}{B}n\log n}$ totali;
(b) qsort permuta \emph{puntatori}: \kw{gli accessi ai suffissi sono sparsi e distruggono la
localit\`a} (niente prefetching). \`E l'esempio da manuale di algoritmo RAM-accettabile ma
I/O-pessimo; le costruzioni serie (skew/DC3, SA-IS) ottengono $\OO{n}$ o $\OO{sort(n)}$.

\textbf{Caso LCP massimo $\le c$.} Ogni confronto termina entro $c+1$ caratteri, cio\`e $\OO{1}$:
\[
\boxed{\ \OO{n\log n}\ \text{tempo}\ }
\]
\kw{Il degrado dipende dalla ripetitivit\`a del testo, misurata dagli LCP.}

\subsection*{\qid{7.6} Suffix tree: la pi\`u lunga sottostringa che occorre $\ge L$ volte}

\begin{domanda}
Dato il suffix tree di $T[1,n]$, come si calcola la sottostringa pi\`u lunga che occorre almeno
$L$ volte? Complessit\`a in funzione di $n$. \emph{(12/06/20)}
\end{domanda}

\textbf{Fatti usati.} (i) ogni sottostringa di $T$ \`e un cammino dalla radice (anche a met\`a
arco); (ii) \kw{il numero di occorrenze di una sottostringa $=$ numero di foglie nel sottoalbero}
del punto che la rappresenta; (iii) i nodi interni sono $\le n-1$.

\textbf{Algoritmo.}
\begin{enumerate}[leftmargin=1.8em,itemsep=1pt]
\item Costruisci il suffix tree (dato, oppure $\OO{n}$ con McCreight/Ukkonen).
\item Con una DFS annota ogni nodo interno $v$ con $\mathit{leaves}(v)$ (somma sui figli, foglie
$=1$): $\OO{n}$.
\item Fra i nodi interni con $\mathit{leaves}(v)\ge L$, prendi quello di \kw{massima string
depth}: la stringa radice$\to v^{*}$ \`e la risposta.
\end{enumerate}

\textbf{Correttezza.} A met\`a arco il sottoalbero (quindi il numero di occorrenze) \`e lo stesso
del nodo sottostante: la risposta ottima pu\`o essere estesa fino al prossimo nodo interno senza
perdere occorrenze, quindi \kw{termina esattamente in un nodo interno}, che l'algoritmo trova.
\textbf{Complessit\`a}: $\boxed{\OO{n}}$, indipendente da $L$.

\subsection*{\qid{7.7} Array LCP in tempo lineare (algoritmo di Kasai)}

\begin{domanda}
Scrivi lo pseudocodice dell'algoritmo lineare (Kasai et al.) che calcola l'array LCP dato $T$ e
$SA$, e dimostra che \`e $\OO{n}$. \emph{(orale 25/07/22, domanda ``da lode''; simulazioni in
11/11/24, 06/03/25, 15/07/25)}
\end{domanda}

Serve l'inverso $\mathit{ISA}$ ($\mathit{ISA}[\mathit{SA}[i]]=i$), calcolabile in $\OO{n}$.

\begin{verbatim}
Kasai(T, SA):
  for i = 1..n: ISA[SA[i]] = i        // inverso di SA
  h = 0
  for i = 1..n:                       // i = posizione NEL TESTO
      if ISA[i] > 1:
           j = SA[ ISA[i] - 1 ]       // suffisso che precede T[i,n] in ordine less.
           while i+h <= n and j+h <= n and T[i+h] == T[j+h]:
                h = h + 1             // estende il prefisso comune
           LCP[ ISA[i] - 1 ] = h
           if h > 0: h = h - 1        // proprieta' chiave (vedi prova)
      else:
           h = 0                      // T[i,n] e' il minimo lessicografico
  return LCP
\end{verbatim}

\textbf{Idea e correttezza.} Si processano i suffissi \kw{in ordine di posizione nel testo}, non
lessicografico, riusando il lavoro: se $T[i,n]$ ha lcp $=h$ col suo predecessore lessicografico
$T[j,n]$, togliendo il primo carattere si ottengono $T[i+1,n]$ e $T[j+1,n]$ con prefisso comune
$h-1$, e $T[j+1,n]$ \`e ancora minore di $T[i+1,n]$. Il predecessore \emph{immediato} di
$T[i+1,n]$ \`e, fra i minori, quello con lcp massimo: quindi \kw{il suo lcp \`e $\ge h-1$} e il
\texttt{while} del passo $i+1$ pu\`o ripartire da $h-1$ senza ricontrollare nulla.

\textbf{Complessit\`a $\OO{n}$ (prova).} Il costo dominante \`e il numero totale di incrementi di
$h$. La variabile $h$: parte da 0, \kw{non supera mai $n$, e viene decrementata al pi\`u $n$
volte} (una per iterazione del for, mai sotto 0). Quindi gli incrementi sono al pi\`u
(decrementi) $+$ (massimo) $\le 2n$:
\[
\boxed{\ \le 2n\ \text{confronti di caratteri in totale}\ \Rightarrow\ \OO{n}\ }
\]
$\blacksquare$ (\`E anche la risposta alla richiesta ``mostra il numero di caratteri
effettivamente confrontati'' delle simulazioni d'esame.)
